% ============================================================
% Shared preamble for all variants of the document.
% Each top-level file (main.tex, main-algebra.tex,
% main-projects.tex) does:
%     \providecommand{\DocPdfTitle}{...}   % optional override
%     % ============================================================
% Shared preamble for all variants of the document.
% Each top-level file (main.tex, main-algebra.tex,
% main-projects.tex) does:
%     \providecommand{\DocPdfTitle}{...}   % optional override
%     % ============================================================
% Shared preamble for all variants of the document.
% Each top-level file (main.tex, main-algebra.tex,
% main-projects.tex) does:
%     \providecommand{\DocPdfTitle}{...}   % optional override
%     % ============================================================
% Shared preamble for all variants of the document.
% Each top-level file (main.tex, main-algebra.tex,
% main-projects.tex) does:
%     \providecommand{\DocPdfTitle}{...}   % optional override
%     \input{preamble}
% then supplies its own title/abstract and \input's the sections.
% ============================================================
\documentclass[11pt]{article}

\usepackage[T1]{fontenc}
\usepackage[utf8]{inputenc}
\usepackage[margin=1.15in]{geometry}

% ---- Math (load before newpx; newpxmath supplies AMS symbols) ----
\usepackage{amsmath,amsthm}

% ---- Typography ----
\usepackage{newpxtext,newpxmath}   % Palatino-family text + matching math
\usepackage{microtype}

% ---- Tables ----
\usepackage{booktabs,array}
\usepackage{float}   % [H] to pin the batteries table at its reference

% ---- Color, graphics, boxes ----
\usepackage[dvipsnames]{xcolor}
\usepackage{graphicx}
\usepackage{tikz}
\usetikzlibrary{arrows.meta,positioning,calc,shapes.geometric,decorations.pathreplacing}
\usepackage[most]{tcolorbox}

\definecolor{linknavy}{RGB}{0,60,120}
\definecolor{citewine}{RGB}{120,20,50}
\definecolor{projfill}{RGB}{245,247,240}
\definecolor{projframe}{RGB}{90,120,70}
\definecolor{msgfill}{RGB}{213,229,245}   % message bits (blue)
\definecolor{lenfill}{RGB}{250,224,188}   % length field (orange)
\definecolor{padfill}{RGB}{198,224,180}   % single pad bit (green)
\definecolor{archgreen}{RGB}{20,110,60}   % local-archive links in bibliography
\definecolor{codegray}{RGB}{120,120,120}

% ---- Code listings (appendix) ----
\usepackage{listings}
\lstset{
  language=Python,
  basicstyle=\ttfamily\footnotesize,
  keywordstyle=\bfseries\color{linknavy},
  commentstyle=\itshape\color{codegray},
  stringstyle=\color{archgreen},
  showstringspaces=false,
  numbers=left,
  numberstyle=\tiny\color{codegray},
  numbersep=8pt,
  breaklines=true,
  keepspaces=true,
  columns=flexible,
  frame=leftline,
  framerule=0.8pt,
  rulecolor=\color{projframe},
  xleftmargin=1.5em,
  aboveskip=1em,
}

% ---- Citations: natbib (author-year) + BibTeX ----
\usepackage[round]{natbib}

% ---- Hyperlinks with SECTION-level back-references ----
% backref=section makes each bibliography entry list the *sections*
% in which it is cited, hyperlinked back into the text.
\providecommand{\DocPdfTitle}{SHA-256 as a Mathematical Object}
\usepackage[colorlinks=true,
            linkcolor=linknavy,
            citecolor=citewine,
            urlcolor=linknavy,
            backref=section,
            pdftitle={\DocPdfTitle},
            pdfauthor={}]{hyperref}
\usepackage{cleveref}

% Last-updated stamp: `make` regenerates lastupdated.tex with the
% current Pacific date/time; fall back to \today if absent.
\InputIfFileExists{lastupdated}{}{}
\providecommand{\lastupdated}{\today}

% The bibliography is thematically clustered: tools/clusterbbl.py
% (run automatically after bibtex via .latexmkrc) inserts cluster
% headers into the .bbl. Suppress natbib's automatic heading; we
% typeset our own.
\renewcommand{\bibsection}{}

% Format back-references as "Cited in §2, §4.1, and §6."
\renewcommand*{\backrefsep}{, \S}
\renewcommand*{\backreftwosep}{ and \S}
\renewcommand*{\backreflastsep}{, and \S}
\renewcommand*{\backref}[1]{}% disable the plain list; use \backrefalt
\renewcommand*{\backrefalt}[4]{%
  \ifcase #1\relax
  \or {\footnotesize\itshape Cited in \S#2.}%
  \else {\footnotesize\itshape Cited in \S#2.}%
  \fi}

% ---- Theorem environments ----
\newtheorem{theorem}{Theorem}[section]
\newtheorem{fact}[theorem]{Fact}
\newtheorem{conjecture}[theorem]{Conjecture}
\theoremstyle{definition}
\newtheorem{definition}[theorem]{Definition}
\theoremstyle{remark}
\newtheorem{remark}[theorem]{Remark}

% ---- Undergraduate project boxes ----
% \begin{project}{Short title}{prerequisites}{time window} ... \end{project}
% Plain (non-\refstepcounter) counter: keeps citations inside boxes
% anchored to the enclosing section, so backref dedup works.
\newcounter{projectctr}
\newtcolorbox{project}[3]{%
  enhanced, breakable,
  code={\stepcounter{projectctr}},
  colback=projfill, colframe=projframe,
  fonttitle=\bfseries\small,
  title={Project \theprojectctr: #1},
  attach boxed title to top left={yshift=-2mm, xshift=4mm},
  boxed title style={colback=projframe},
  after upper={\par\smallskip\footnotesize\textit{Prerequisites:} #2\quad
               \textit{Window:} #3},
  left=2mm, right=2mm, top=3mm, bottom=2mm
}

% ---- Reading-spine box for variant front matter ----
% \begin{readingspine}{title} ... \end{readingspine}
\newtcolorbox{readingspine}[1]{%
  enhanced, breakable,
  colback=projfill, colframe=projframe,
  fonttitle=\bfseries\small,
  title={#1},
  attach boxed title to top left={yshift=-2mm, xshift=4mm},
  boxed title style={colback=projframe},
  left=2mm, right=2mm, top=3mm, bottom=2mm
}

% Local archived copy of a cited work: visually distinct (green) and
% clickable — opens refs/<file> relative to this PDF's location.
\newcommand{\localpdf}[1]{\href{./refs/#1}{\textcolor{archgreen}{$\blacktriangleright$\,\texttt{refs/#1}}}}

% ---- Notation macros ----
\newcommand{\Ztwo}{\mathbb{Z}/2^{32}\mathbb{Z}}
\newcommand{\Ftwo}{\mathbb{F}_2}
\newcommand{\State}{\mathcal{S}}
\newcommand{\shaxor}{\oplus}
\newcommand{\shaplus}{\boxplus}
\newcommand{\SHA}{\textsf{SHA-256}}
\newcommand{\bitstrings}{\{0,1\}}

% then supplies its own title/abstract and \input's the sections.
% ============================================================
\documentclass[11pt]{article}

\usepackage[T1]{fontenc}
\usepackage[utf8]{inputenc}
\usepackage[margin=1.15in]{geometry}

% ---- Math (load before newpx; newpxmath supplies AMS symbols) ----
\usepackage{amsmath,amsthm}

% ---- Typography ----
\usepackage{newpxtext,newpxmath}   % Palatino-family text + matching math
\usepackage{microtype}

% ---- Tables ----
\usepackage{booktabs,array}
\usepackage{float}   % [H] to pin the batteries table at its reference

% ---- Color, graphics, boxes ----
\usepackage[dvipsnames]{xcolor}
\usepackage{graphicx}
\usepackage{tikz}
\usetikzlibrary{arrows.meta,positioning,calc,shapes.geometric,decorations.pathreplacing}
\usepackage[most]{tcolorbox}

\definecolor{linknavy}{RGB}{0,60,120}
\definecolor{citewine}{RGB}{120,20,50}
\definecolor{projfill}{RGB}{245,247,240}
\definecolor{projframe}{RGB}{90,120,70}
\definecolor{msgfill}{RGB}{213,229,245}   % message bits (blue)
\definecolor{lenfill}{RGB}{250,224,188}   % length field (orange)
\definecolor{padfill}{RGB}{198,224,180}   % single pad bit (green)
\definecolor{archgreen}{RGB}{20,110,60}   % local-archive links in bibliography
\definecolor{codegray}{RGB}{120,120,120}

% ---- Code listings (appendix) ----
\usepackage{listings}
\lstset{
  language=Python,
  basicstyle=\ttfamily\footnotesize,
  keywordstyle=\bfseries\color{linknavy},
  commentstyle=\itshape\color{codegray},
  stringstyle=\color{archgreen},
  showstringspaces=false,
  numbers=left,
  numberstyle=\tiny\color{codegray},
  numbersep=8pt,
  breaklines=true,
  keepspaces=true,
  columns=flexible,
  frame=leftline,
  framerule=0.8pt,
  rulecolor=\color{projframe},
  xleftmargin=1.5em,
  aboveskip=1em,
}

% ---- Citations: natbib (author-year) + BibTeX ----
\usepackage[round]{natbib}

% ---- Hyperlinks with SECTION-level back-references ----
% backref=section makes each bibliography entry list the *sections*
% in which it is cited, hyperlinked back into the text.
\providecommand{\DocPdfTitle}{SHA-256 as a Mathematical Object}
\usepackage[colorlinks=true,
            linkcolor=linknavy,
            citecolor=citewine,
            urlcolor=linknavy,
            backref=section,
            pdftitle={\DocPdfTitle},
            pdfauthor={}]{hyperref}
\usepackage{cleveref}

% Last-updated stamp: `make` regenerates lastupdated.tex with the
% current Pacific date/time; fall back to \today if absent.
\InputIfFileExists{lastupdated}{}{}
\providecommand{\lastupdated}{\today}

% The bibliography is thematically clustered: tools/clusterbbl.py
% (run automatically after bibtex via .latexmkrc) inserts cluster
% headers into the .bbl. Suppress natbib's automatic heading; we
% typeset our own.
\renewcommand{\bibsection}{}

% Format back-references as "Cited in §2, §4.1, and §6."
\renewcommand*{\backrefsep}{, \S}
\renewcommand*{\backreftwosep}{ and \S}
\renewcommand*{\backreflastsep}{, and \S}
\renewcommand*{\backref}[1]{}% disable the plain list; use \backrefalt
\renewcommand*{\backrefalt}[4]{%
  \ifcase #1\relax
  \or {\footnotesize\itshape Cited in \S#2.}%
  \else {\footnotesize\itshape Cited in \S#2.}%
  \fi}

% ---- Theorem environments ----
\newtheorem{theorem}{Theorem}[section]
\newtheorem{fact}[theorem]{Fact}
\newtheorem{conjecture}[theorem]{Conjecture}
\theoremstyle{definition}
\newtheorem{definition}[theorem]{Definition}
\theoremstyle{remark}
\newtheorem{remark}[theorem]{Remark}

% ---- Undergraduate project boxes ----
% \begin{project}{Short title}{prerequisites}{time window} ... \end{project}
% Plain (non-\refstepcounter) counter: keeps citations inside boxes
% anchored to the enclosing section, so backref dedup works.
\newcounter{projectctr}
\newtcolorbox{project}[3]{%
  enhanced, breakable,
  code={\stepcounter{projectctr}},
  colback=projfill, colframe=projframe,
  fonttitle=\bfseries\small,
  title={Project \theprojectctr: #1},
  attach boxed title to top left={yshift=-2mm, xshift=4mm},
  boxed title style={colback=projframe},
  after upper={\par\smallskip\footnotesize\textit{Prerequisites:} #2\quad
               \textit{Window:} #3},
  left=2mm, right=2mm, top=3mm, bottom=2mm
}

% ---- Reading-spine box for variant front matter ----
% \begin{readingspine}{title} ... \end{readingspine}
\newtcolorbox{readingspine}[1]{%
  enhanced, breakable,
  colback=projfill, colframe=projframe,
  fonttitle=\bfseries\small,
  title={#1},
  attach boxed title to top left={yshift=-2mm, xshift=4mm},
  boxed title style={colback=projframe},
  left=2mm, right=2mm, top=3mm, bottom=2mm
}

% Local archived copy of a cited work: visually distinct (green) and
% clickable — opens refs/<file> relative to this PDF's location.
\newcommand{\localpdf}[1]{\href{./refs/#1}{\textcolor{archgreen}{$\blacktriangleright$\,\texttt{refs/#1}}}}

% ---- Notation macros ----
\newcommand{\Ztwo}{\mathbb{Z}/2^{32}\mathbb{Z}}
\newcommand{\Ftwo}{\mathbb{F}_2}
\newcommand{\State}{\mathcal{S}}
\newcommand{\shaxor}{\oplus}
\newcommand{\shaplus}{\boxplus}
\newcommand{\SHA}{\textsf{SHA-256}}
\newcommand{\bitstrings}{\{0,1\}}

% then supplies its own title/abstract and \input's the sections.
% ============================================================
\documentclass[11pt]{article}

\usepackage[T1]{fontenc}
\usepackage[utf8]{inputenc}
\usepackage[margin=1.15in]{geometry}

% ---- Math (load before newpx; newpxmath supplies AMS symbols) ----
\usepackage{amsmath,amsthm}

% ---- Typography ----
\usepackage{newpxtext,newpxmath}   % Palatino-family text + matching math
\usepackage{microtype}

% ---- Tables ----
\usepackage{booktabs,array}
\usepackage{float}   % [H] to pin the batteries table at its reference

% ---- Color, graphics, boxes ----
\usepackage[dvipsnames]{xcolor}
\usepackage{graphicx}
\usepackage{tikz}
\usetikzlibrary{arrows.meta,positioning,calc,shapes.geometric,decorations.pathreplacing}
\usepackage[most]{tcolorbox}

\definecolor{linknavy}{RGB}{0,60,120}
\definecolor{citewine}{RGB}{120,20,50}
\definecolor{projfill}{RGB}{245,247,240}
\definecolor{projframe}{RGB}{90,120,70}
\definecolor{msgfill}{RGB}{213,229,245}   % message bits (blue)
\definecolor{lenfill}{RGB}{250,224,188}   % length field (orange)
\definecolor{padfill}{RGB}{198,224,180}   % single pad bit (green)
\definecolor{archgreen}{RGB}{20,110,60}   % local-archive links in bibliography
\definecolor{codegray}{RGB}{120,120,120}

% ---- Code listings (appendix) ----
\usepackage{listings}
\lstset{
  language=Python,
  basicstyle=\ttfamily\footnotesize,
  keywordstyle=\bfseries\color{linknavy},
  commentstyle=\itshape\color{codegray},
  stringstyle=\color{archgreen},
  showstringspaces=false,
  numbers=left,
  numberstyle=\tiny\color{codegray},
  numbersep=8pt,
  breaklines=true,
  keepspaces=true,
  columns=flexible,
  frame=leftline,
  framerule=0.8pt,
  rulecolor=\color{projframe},
  xleftmargin=1.5em,
  aboveskip=1em,
}

% ---- Citations: natbib (author-year) + BibTeX ----
\usepackage[round]{natbib}

% ---- Hyperlinks with SECTION-level back-references ----
% backref=section makes each bibliography entry list the *sections*
% in which it is cited, hyperlinked back into the text.
\providecommand{\DocPdfTitle}{SHA-256 as a Mathematical Object}
\usepackage[colorlinks=true,
            linkcolor=linknavy,
            citecolor=citewine,
            urlcolor=linknavy,
            backref=section,
            pdftitle={\DocPdfTitle},
            pdfauthor={}]{hyperref}
\usepackage{cleveref}

% Last-updated stamp: `make` regenerates lastupdated.tex with the
% current Pacific date/time; fall back to \today if absent.
\InputIfFileExists{lastupdated}{}{}
\providecommand{\lastupdated}{\today}

% The bibliography is thematically clustered: tools/clusterbbl.py
% (run automatically after bibtex via .latexmkrc) inserts cluster
% headers into the .bbl. Suppress natbib's automatic heading; we
% typeset our own.
\renewcommand{\bibsection}{}

% Format back-references as "Cited in §2, §4.1, and §6."
\renewcommand*{\backrefsep}{, \S}
\renewcommand*{\backreftwosep}{ and \S}
\renewcommand*{\backreflastsep}{, and \S}
\renewcommand*{\backref}[1]{}% disable the plain list; use \backrefalt
\renewcommand*{\backrefalt}[4]{%
  \ifcase #1\relax
  \or {\footnotesize\itshape Cited in \S#2.}%
  \else {\footnotesize\itshape Cited in \S#2.}%
  \fi}

% ---- Theorem environments ----
\newtheorem{theorem}{Theorem}[section]
\newtheorem{fact}[theorem]{Fact}
\newtheorem{conjecture}[theorem]{Conjecture}
\theoremstyle{definition}
\newtheorem{definition}[theorem]{Definition}
\theoremstyle{remark}
\newtheorem{remark}[theorem]{Remark}

% ---- Undergraduate project boxes ----
% \begin{project}{Short title}{prerequisites}{time window} ... \end{project}
% Plain (non-\refstepcounter) counter: keeps citations inside boxes
% anchored to the enclosing section, so backref dedup works.
\newcounter{projectctr}
\newtcolorbox{project}[3]{%
  enhanced, breakable,
  code={\stepcounter{projectctr}},
  colback=projfill, colframe=projframe,
  fonttitle=\bfseries\small,
  title={Project \theprojectctr: #1},
  attach boxed title to top left={yshift=-2mm, xshift=4mm},
  boxed title style={colback=projframe},
  after upper={\par\smallskip\footnotesize\textit{Prerequisites:} #2\quad
               \textit{Window:} #3},
  left=2mm, right=2mm, top=3mm, bottom=2mm
}

% ---- Reading-spine box for variant front matter ----
% \begin{readingspine}{title} ... \end{readingspine}
\newtcolorbox{readingspine}[1]{%
  enhanced, breakable,
  colback=projfill, colframe=projframe,
  fonttitle=\bfseries\small,
  title={#1},
  attach boxed title to top left={yshift=-2mm, xshift=4mm},
  boxed title style={colback=projframe},
  left=2mm, right=2mm, top=3mm, bottom=2mm
}

% Local archived copy of a cited work: visually distinct (green) and
% clickable — opens refs/<file> relative to this PDF's location.
\newcommand{\localpdf}[1]{\href{./refs/#1}{\textcolor{archgreen}{$\blacktriangleright$\,\texttt{refs/#1}}}}

% ---- Notation macros ----
\newcommand{\Ztwo}{\mathbb{Z}/2^{32}\mathbb{Z}}
\newcommand{\Ftwo}{\mathbb{F}_2}
\newcommand{\State}{\mathcal{S}}
\newcommand{\shaxor}{\oplus}
\newcommand{\shaplus}{\boxplus}
\newcommand{\SHA}{\textsf{SHA-256}}
\newcommand{\bitstrings}{\{0,1\}}

% then supplies its own title/abstract and \input's the sections.
% ============================================================
\documentclass[11pt]{article}

\usepackage[T1]{fontenc}
\usepackage[utf8]{inputenc}
\usepackage[margin=1.15in]{geometry}

% ---- Math (load before newpx; newpxmath supplies AMS symbols) ----
\usepackage{amsmath,amsthm}

% ---- Typography ----
\usepackage{newpxtext,newpxmath}   % Palatino-family text + matching math
\usepackage{microtype}

% ---- Tables ----
\usepackage{booktabs,array}
\usepackage{float}   % [H] to pin the batteries table at its reference

% ---- Color, graphics, boxes ----
\usepackage[dvipsnames]{xcolor}
\usepackage{graphicx}
\usepackage{tikz}
\usetikzlibrary{arrows.meta,positioning,calc,shapes.geometric,decorations.pathreplacing}
\usepackage[most]{tcolorbox}

\definecolor{linknavy}{RGB}{0,60,120}
\definecolor{citewine}{RGB}{120,20,50}
\definecolor{projfill}{RGB}{245,247,240}
\definecolor{projframe}{RGB}{90,120,70}
\definecolor{msgfill}{RGB}{213,229,245}   % message bits (blue)
\definecolor{lenfill}{RGB}{250,224,188}   % length field (orange)
\definecolor{padfill}{RGB}{198,224,180}   % single pad bit (green)
\definecolor{archgreen}{RGB}{20,110,60}   % local-archive links in bibliography
\definecolor{codegray}{RGB}{120,120,120}

% ---- Code listings (appendix) ----
\usepackage{listings}
\lstset{
  language=Python,
  basicstyle=\ttfamily\footnotesize,
  keywordstyle=\bfseries\color{linknavy},
  commentstyle=\itshape\color{codegray},
  stringstyle=\color{archgreen},
  showstringspaces=false,
  numbers=left,
  numberstyle=\tiny\color{codegray},
  numbersep=8pt,
  breaklines=true,
  keepspaces=true,
  columns=flexible,
  frame=leftline,
  framerule=0.8pt,
  rulecolor=\color{projframe},
  xleftmargin=1.5em,
  aboveskip=1em,
}

% ---- Citations: natbib (author-year) + BibTeX ----
\usepackage[round]{natbib}

% ---- Hyperlinks with SECTION-level back-references ----
% backref=section makes each bibliography entry list the *sections*
% in which it is cited, hyperlinked back into the text.
\providecommand{\DocPdfTitle}{SHA-256 as a Mathematical Object}
\usepackage[colorlinks=true,
            linkcolor=linknavy,
            citecolor=citewine,
            urlcolor=linknavy,
            backref=section,
            pdftitle={\DocPdfTitle},
            pdfauthor={}]{hyperref}
\usepackage{cleveref}

% Last-updated stamp: `make` regenerates lastupdated.tex with the
% current Pacific date/time; fall back to \today if absent.
\InputIfFileExists{lastupdated}{}{}
\providecommand{\lastupdated}{\today}

% The bibliography is thematically clustered: tools/clusterbbl.py
% (run automatically after bibtex via .latexmkrc) inserts cluster
% headers into the .bbl. Suppress natbib's automatic heading; we
% typeset our own.
\renewcommand{\bibsection}{}

% Format back-references as "Cited in §2, §4.1, and §6."
\renewcommand*{\backrefsep}{, \S}
\renewcommand*{\backreftwosep}{ and \S}
\renewcommand*{\backreflastsep}{, and \S}
\renewcommand*{\backref}[1]{}% disable the plain list; use \backrefalt
\renewcommand*{\backrefalt}[4]{%
  \ifcase #1\relax
  \or {\footnotesize\itshape Cited in \S#2.}%
  \else {\footnotesize\itshape Cited in \S#2.}%
  \fi}

% ---- Theorem environments ----
\newtheorem{theorem}{Theorem}[section]
\newtheorem{fact}[theorem]{Fact}
\newtheorem{conjecture}[theorem]{Conjecture}
\theoremstyle{definition}
\newtheorem{definition}[theorem]{Definition}
\theoremstyle{remark}
\newtheorem{remark}[theorem]{Remark}

% ---- Undergraduate project boxes ----
% \begin{project}{Short title}{prerequisites}{time window} ... \end{project}
% Plain (non-\refstepcounter) counter: keeps citations inside boxes
% anchored to the enclosing section, so backref dedup works.
\newcounter{projectctr}
\newtcolorbox{project}[3]{%
  enhanced, breakable,
  code={\stepcounter{projectctr}},
  colback=projfill, colframe=projframe,
  fonttitle=\bfseries\small,
  title={Project \theprojectctr: #1},
  attach boxed title to top left={yshift=-2mm, xshift=4mm},
  boxed title style={colback=projframe},
  after upper={\par\smallskip\footnotesize\textit{Prerequisites:} #2\quad
               \textit{Window:} #3},
  left=2mm, right=2mm, top=3mm, bottom=2mm
}

% ---- Reading-spine box for variant front matter ----
% \begin{readingspine}{title} ... \end{readingspine}
\newtcolorbox{readingspine}[1]{%
  enhanced, breakable,
  colback=projfill, colframe=projframe,
  fonttitle=\bfseries\small,
  title={#1},
  attach boxed title to top left={yshift=-2mm, xshift=4mm},
  boxed title style={colback=projframe},
  left=2mm, right=2mm, top=3mm, bottom=2mm
}

% Local archived copy of a cited work: visually distinct (green) and
% clickable — opens refs/<file> relative to this PDF's location.
\newcommand{\localpdf}[1]{\href{./refs/#1}{\textcolor{archgreen}{$\blacktriangleright$\,\texttt{refs/#1}}}}

% ---- Notation macros ----
\newcommand{\Ztwo}{\mathbb{Z}/2^{32}\mathbb{Z}}
\newcommand{\Ftwo}{\mathbb{F}_2}
\newcommand{\State}{\mathcal{S}}
\newcommand{\shaxor}{\oplus}
\newcommand{\shaplus}{\boxplus}
\newcommand{\SHA}{\textsf{SHA-256}}
\newcommand{\bitstrings}{\{0,1\}}
